\exnumber{01}
\extitle{UI IV - Ex 2 - Design system library}
\exfiles{ex01/Figma file, named Name\_FirstName\_DeliverableName\_Date\_VersionNumber}
\exauthorize{Those named in the exercise}
\makeheaderfiles

\textbf{Instructions:}

Based on the moodboard you designed in Ex 1, create a design system with at least 30 elements, including atoms, molecules and organisms, and include multistates.

\textbf{Deliverables:}

\begin{itemize}
\item A Figma file with a library including:
\begin{itemize}
\item At least graphic design 30 elements
\item At least 5 atoms
\item At least 5 molecules
\item At least 5 organisms
\item At least 10 multistate variations
\end{itemize}
\item Upload the Figma file on Github, labeled as follows: "Name\_FirstName\_DeliverableName\_Date\_VersionNumber". For example, it can be labeled as "Doe\_John\_MidFiPrototype\_05242024\_V1".
\item Add a title within the file.
\item Pay attention to the global aspect of the file. It must be clear, simple and easy to read. You can get inspiration on canva.com but don't overload your design with too much details!
\end{itemize}

\textbf{Tips:}

\begin{itemize}
\item Start by asking yourself what key elements may appear in an e-commerce app.
\item Feel free to get inspiration from existing design systems. Browse Dribble and the Figma community!
\end{itemize}

\textbf{Resources:}

\begin{itemize}
\item FIGMA FILE (to duplicate) \url{https://www.figma.com/file/U26mctkrOqxiUeAlNHzSxW/UI-IV---Ex-3---E-COMMERCE%3A-Design-screens!?node-id=0%3A1}
\item Introduction to design systems \url{https://www.youtube.com/watch?v=xEuBCUngJ_U}
\item Create a design system - Figma course \url{https://www.youtube.com/watch?v=RYDiDpW2VkM}
\item Building a design system library - Lyft \url{https://medium.com/tap-to-dismiss/building-a-design-system-library-532ef2492811}
\item Examples of UI Buttons library \url{https://www.pinterest.co.kr/pin/663014376380850291/}
\end{itemize}

